\chapter{Monodromic generation by survival of
\texorpdfstring{\(\pi\), \(e\) and \(\varphi\)}{pi, e and phi}}
\label{chap:generacion-numeros-rev6}
\label{chap:biografias-numeros-rev6}
\index{HMT!number generation by survival}
\index{pi@\(\pi\)!regional closure}
\index{e@\(e\)!propagation and memory}
\index{phi@\(\varphi\)!self-scaling}
\index{nonadic return}\index{mirror pi/e@mirror \(\pi/e\)}

Generation does not begin with a prescribed decimal expansion. It begins in
the complete TPK state: visible word, sheet, orientation, quotient,
carry, Hensel state, ternary and decimal cylinders, joint signature,
boundary, phase, and dodecaphasic register. The EF--\(G_9\) relation prolongs that
state and eliminates a continuation only when it ceases to survive some
future boundary. The decimal reader subsequently acts on the section that has
been determined.

The causal chain is
\[
\Omega_{\rm TPK}
\longrightarrow 468\,\mathcal U_6
\longrightarrow 243\,\mathcal W_6
\longrightarrow
\bigl(w_\pi,w_e,w_\varphi\bigr)
\longrightarrow
\bigl(\mathcal R_{\pi,n},\mathcal R_{e,n},
      \mathcal R_{\varphi,n}\bigr)_{n\geq0}
\longrightarrow
\varprojlim_n\mathcal R_{\chi,n}
\xrightarrow{\ L_{\rm arch}\ }
\chi .
\]
The ambient space \(\mathbb F_3^6\) has \(729\) words, but the effective
image contains \(243\); the two objects are not identified. Within that
image, internal selection fixes
\[
w_\pi=010211,\qquad
w_e=201101,\qquad
w_\varphi=121200
\]
Before consulting figures of \(\pi,e,\varphi\).

\section{Generative regions and global section}

\begin{definition}[Generative region of a channel]
\label{def:region-generativa-rev6}
For \(\chi\in\{\pi,e,\varphi\}\), let \(\mathcal R_{\chi,n}\) be the
finite set of complete states of depth \(n\) that start from
\(w_\chi\) and satisfy at each transition the relation EF--\(G_9\), the
exact intersection of cylinders, and the joint boundary signature. If
\(\tau_{n+1,n}\) eliminates only the last layer, we define
\[
\mathcal R_\chi
=\varprojlim_n
\bigl(\mathcal R_{\chi,n},\tau_{n+1,n}\bigr).
\]
The published value is \(L_{\rm arch}(x_\chi)\), where
\(x_\chi\in\mathcal R_\chi\); the reader is not part of the state that
selects it.
\end{definition}

\begin{theorem}[Monodromy Prolongation Principle in the Full TPK]
\label{thm:generacion-supervivencia-pi-e-phi-rev6}
Suppose that, for a channel
\(\chi\in\{\pi,e,\varphi\}\), each set
\(\operatorname{Surv}_{\chi,n}\) is nonempty and unitary, and that its elements
are compatible with truncation. Then there exists a unique section
\[
x_\chi=(x_{\chi,n})_{n\geq0}\in\mathcal R_\chi
\]
The decimal cylinders associated with their truncations are nested, and their
diameters tend to zero. If the channel reader identifies their intervals
with the declared closure, propagation, or self-scaling, the intersection
publishes, respectively, \(\pi\), \(e\), or \(\varphi\).
\end{theorem}

\begin{proof}
In the full profile, stable survival at level \(n\) is
\[
\operatorname{Surv}_{\chi,n}
=\bigcap_{h\geq1}
\left\{s:\text{\(s\) admits a compatible prolongation across \(h\) boundaries}\right\}.
\]
By hypothesis, the surviving state at each level truncates to that at the
preceding level. The levels form a finitely branching, nonempty tree
at every depth. The Koenig lemma provides a global branch;
levelwise unitarity precludes a second section.

For a ternary prefix of length \(L\), interpreted as \(N\), and a
prefix of \(m\) triads, interpreted as \(D\), the compatibility conditions are
\[
I_{N,L}=\left[\frac{N}{3^L},\frac{N+1}{3^L}\right),
\qquad
J_{D,m}=\left[\frac{D}{1000^m},\frac{D+1}{1000^m}\right),
\qquad
I_{N,L}\cap J_{D,m}\ne\varnothing .
\]
The prolongation conserves the intersection and makes both
diameters tend to zero. The total intersection contains a unique real. The declared
reader identifies that real within its channel.
\end{proof}

The reproducible execution does not replace the theorem's hypotheses with an
implicit assertion. It verifies them at each of the (3501) levels
needed to print ten thousand digits: there the partition \(729\to1\), the
compatibility of truncations, and the change of memory after return are
checked exhaustively.

\section{Survival filtration, nonadic monodromy, and \texorpdfstring{\(\pi/e\)}{pi/e} symmetry}

For
\[
B_n=(u_n^\pi,u_n^e,u_n^\varphi),\qquad
q_n=u_n^\pi+u_n^e+u_n^\varphi,\qquad
a_n=u_n^\pi+u_n^e-u_n^\varphi,
\]
the joint signature determines
\[
u_n^\varphi=2(q_n-a_n),\qquad
u_n^\pi+u_n^e=q_n-u_n^\varphi
\quad\text{in }\mathbb F_3^6.
\]
Thus, \(\varphi\) is the fixed self-scaling axis, and the only local ambiguity
is the mirror allocation between \(\pi\) and \(e\).

The phase is \(g(n)=n\bmod9\), with \(0\equiv9\). After \(w_{30}\), the
first turn traverses
\[
5,6,7,8,9,1,2,3,4.
\]
They are not nine independent filters, but the nine phases of a single
connection: pruning, synchronization, mirror orientation, reopening, and
confirmation by survival.

In the ninth phase,
\[
B_9=(100100\mid020112\mid010122)
\]
has three locally compatible outputs. The three conserve
\(u_{10}^\varphi=112002\) and
\(u_{10}^\pi+u_{10}^e=210110\); all survive one boundary, but only
\[
\boxed{B_{10}=(101222\mid112221\mid112002)}
\]
the second survives. The future breaks the local mirror. Moreover,
\[
g(n+9)=g(n),\qquad x_{\chi,n+9}\ne x_{\chi,n}
\]
in the returns: the next stage returns to phase \(1\), not to the
preceding state. The quotient, carry, and boundary turn the circuit into a
helix of memory. This return action with nontrivial transport is the
nonadic monodromy of the TPK.

\section{A reproducible emission of twelve triads}

Thirteen blocks, equivalent to \(78\) trits, suffice for the reader to
\[
D=\left\lfloor\frac{1000^{12}N}{3^{78}}\right\rfloor
\]
publish the following segment directly from the ternary branches:
\[
\begin{array}{c|rrrrrrrrrrrr}
\pi&141&592&653&589&793&238&462&643&383&279&502&884\\
e&718&281&828&459&045&235&360&287&471&352&662&497\\
\varphi&618&033&988&749&894&848&204&586&834&365&638&117.
\end{array}
\]
This table does not feed the transition: it exhibits its output. The iteration of the
same connection, with the return \(9\to1\), produces any requested finite prefix
under the hypotheses of the
\cref{thm:generacion-supervivencia-pi-e-phi-rev6}; that theorem specifies the
prolongation to the inverse limit.

\section{Unique closure of \texorpdfstring{\(\pi\)}{pi} and classification of its \(5\) regional histories}

Let
\[
 \Psi:\mathcal U_6\longrightarrow\mathbb F_3^6,
 \qquad
 \Psi(U_6)=(-U_6)\bmod3
\]
the coordinate projection of the catalog. The fiber
\[
 \mathscr R_\pi^{\rm micro}
 :=\Psi^{-1}(010211)
\]
contains exactly the five vectors
\[
\begin{aligned}
U_\pi^\perp&=(501,614,498,169,272,272),\\
U_{\pi,1}^{++}&=(810,923,870,169,272,272),\\
U_{\pi,2}^{++}&=(810,983,810,169,272,272),\\
U_{\pi,3}^{++}&=(870,923,810,169,272,272),\\
U_\pi^{--}&=(870,923,810,418,674,521).
\end{aligned}
\]
Their multiplicative signatures are
\[
555555,\qquad339555,\qquad393555,\qquad933555,\qquad933717,
\]
and their multiplicities are
\[
432,\qquad144,\qquad144,\qquad144,\qquad144.
\]
Therefore,
\[
\boxed{
5=1_\perp+3_{++}+1_{--},
\qquad
1008=432+4\cdot144.}
\]
The first identity classifies orientation functions; the second counts
seed routes. They are not two ways of counting five copies of the same object.

The visible quotient
\[
 \mathscr R_\pi^{w_6}
 =\Psi(\mathscr R_\pi^{\rm micro})
 =\{010211\}
\]
has a single element because it forgets regional provenance. Equality of
words does not identify the five states: the first four regions
preserve the tail \(169|272|272\), whereas the negative region effects a
mutated return. The transverse region, of multiplicity \(432\), likewise cannot be
reduced to one of the four branches of multiplicity \(144\).

The closure function is also seen in the multigraph of blocks. Each of
the five regions belongs to a minimal closed walk of length eight that
also traverses the anchors \(U_e\), \(U_\varphi\), and \(U_{\rm APP}\).
The enumeration and the proof of minimality reside in
\cref{chap:cinco-ciclos-pi}; the regional list and its multiplicities, in
\cref{chap:regiones}.

\begin{proposition}[Closure Package for \(\pi\)]
\label{prop:paquete-clausura-pi-rev6}
The datum
\[
 \mathfrak C_\pi=
 \bigl(
 \mathscr R_\pi^{\rm micro},
 w_\pi,
 \rho_\pi,
 m_\pi,
 \ell_{\rm cl}
 \bigr),
\]
where
\[
\rho_\pi=(\perp,++ ,++ ,++ ,--),\qquad
m_\pi=(432,144,144,144,144),\qquad
\ell_{\rm cl}=8,
\]
preserves exactly the functional information lost by the quotient
\(\mathscr R_\pi^{w_6}\). Consequently, the function of \(\pi\) in this
generation is to close five regional histories into a single evaluation,
not to repeat the same constant five times.
\end{proposition}

\begin{proof}
The fiber, roles, and multiplicities are those of
\cref{chap:regiones}. The \cref{chap:cinco-ciclos-pi} proves that each
region enters a minimal closure of length eight with the three declared
anchors. The projection \(\Psi\) sends the five regions to the same
word, so that the quotient preserves the visible value and forgets
\(\rho_\pi,m_\pi,\ell_{\rm cl}\).
\end{proof}

\section{\texorpdfstring{\(e\)}{e}: visible propagation and differentiated quotient memory}

The fiber of \(w_e=201101\) consists of an emission \(U_e\) with multiplicity
\(144\). Its functional feature is not an infinite periodicity, but rather a
local repetition exhaustively distinguished within the reader of the
\(243\) words. The \cref{thm:cuadrado-local-e} proves that \(w_e\) is the
only word whose twelve-digit reading contains a square of length
four beginning at the second digit:
\[
718281828459
=7\,\underbrace{1828\,1828}_{\text{local square}}\,459.
\]
Juxtaposition of blocks is decimal concatenation; it is not the arithmetic
product \(7\cdot1828^2\cdot459\). The catalog contains another square of
length four,
\[
575157518472
=\underbrace{5751\,5751}_{\text{local square}}\,8472,
\]
but begins at the first digit and, therefore, does not satisfy the same
positional predicate.

The complete census of squares of lengths \(1,\ldots,6\) is
\[
\begin{array}{crrrrrr}
\ell&1&2&3&4&5&6\\
\text{occurrences}&240&21&3&2&0&0\\
\text{words that contain them}&153&19&3&2&0&0.
\end{array}
\]
The repetition \(1828\,1828\) does not begin a period either: a third copy
would require a \(1\) as the next digit, but the effective reading contains
a \(4\).

The dynamics of elevation show why the episode does express
propagation:
\[
201101\mid121221\mid102011\mid012222\mid102011.
\]
The third and fifth visible blocks coincide, but the states that
emit them are distinct. Their prestates modulo \(27\) are
\[
(25,11,7,17,8,16),
\qquad
(7,8,4,23,26,7),
\]
and their subsequent Hensel quotients are
\[
(6,13,9,2,13,1),
\qquad
(15,0,3,19,23,25).
\]
There is return of the projection \(102011\), not return of the enriched state.
The output can propagate and reappear because the quotient transports the
genealogy that the visible word has omitted.

\begin{proposition}[Functional content of \(e\)]
\label{prop:propagacion-e-rev6}
The exact generation of \(e\) combines three compatible facts:
\begin{enumerate}
\item the selection internal of \(w_e=201101\);
\item the positional uniqueness of the local square \(1828\,1828\) in the
fixed catalogue;
\item the reappearance of \(102011\) under different memory states.
\end{enumerate}
Its function is propagation with memory renewal. It is not periodicity of
the expansion nor repetition of the complete state.
\end{proposition}

\begin{proof}
The first point is the case \(e\) of the
\cref{thm:selector-orbital-constantes}; the second and its census are the
\cref{thm:cuadrado-local-e}. The inequality of the two prestates and of the
two Hensel quotients proves the third.
\end{proof}

\section{Autoscaling Axis and Mirror Fiber}

The joint signature of the three generations preserves a linear decomposition
that will subsequently be used by the scale reader. For
\[
 B_k=(u_k^\pi,u_k^e,u_k^\varphi)
 \in(\mathbb F_3^6)^3
\]
we define, coordinate by coordinate,
\[
 q_k=u_k^\pi+u_k^e+u_k^\varphi,
 \qquad
 a_k=u_k^\pi+u_k^e-u_k^\varphi.
\]

\begin{theorem}[Axial decomposition of the boundary signature]
\label{thm:eje-espejo-nonadico}
The pair \((q_k,a_k)\) uniquely determines the self-scaling channel and the
aggregate of the other two channels:
\[
 u_k^\varphi=2(q_k-a_k),
 \qquad
 u_k^\pi+u_k^e=q_k-u_k^\varphi.
\]
Consequently, a fiber of
\[
 q_{\rm ax}(u^\pi,u^e,u^\varphi)
 =(u^\pi+u^e,u^\varphi)
\]
preserves ambiguity only in the internal allocation between \(u^\pi\) and \(u^e\).
\end{theorem}

\begin{proof}
Subtracting the two defining equations gives
\(q_k-a_k=2u_k^\varphi\). Since the inverse of \(2\) in \(\mathbb F_3\) is
again \(2\),
\[
u_k^\varphi=2(q_k-a_k).
\]
Upon subtracting this vector from \(q_k\) there remains
\(u_k^\pi+u_k^e\). The operation fixes the axis \(\varphi\), but does not separate
the two summands of the fiber \(\pi/e\); that is exactly the residual
ambiguity.
\end{proof}

\def\HMTAxialTheoremDefined{1}

The expression \emph{mirror fiber} names this exact decomposition in
\(\mathbb F_3^6\). It does not posit an analytic identity between the real
constants: it indicates which component is fixed by the signature and where an
internal choice persists.

\section{\texorpdfstring{\(\varphi\)}{Phi}: the autoescala inducida by TPK}

The word \(w_\varphi=121200\) possesses a single emission of multiplicity
\(144\), but its function is not deduced from that cardinality. The TPK calendar
produces the block of modes
\[
(\Sigma,\Pi,\Pi).
\]
Reading \(\Sigma\) as the areal leaf and \(\Pi\) as the volumetric leaf, the
three consecutive pairs of the cycle force the incidence
\[
F_{\rm av}
=
\begin{pmatrix}
0&1\\
1&1
\end{pmatrix},
\]
as \cref{thm:descenso-necesario-fav} proves. The golden ratio is not
introduced to choose this matrix: it appears afterward as its Perron root.
Indeed,
\[
\det(tI-F_{\rm av})=t^2-t-1,
\qquad
F_{\rm av}\binom{1}{\varphi}
=\varphi\binom{1}{\varphi}.
\]

For an admissible path \(\gamma:i\to j\) of length \(n\), with
\(r=(1,\varphi)\), the Parry character
\[
\chi_P(\gamma)=\varphi^n\frac{r_i}{r_j}
\]
is multiplicative under concatenation. On a closed path,
\(\chi_P(\gamma)=\varphi^n\). The measure of maximal entropy of the memory-one
hull satisfies, by \cref{cor:parry-envolvente-fav},
\[
\frac{\mu_{\rm ar}}{\mu_{\rm vol}}=\varphi^{-2}.
\]

Second-order memory also preserves the contracted direction. In the basis \((x^2,2xy,y^2)\),
\[
\operatorname{Sym}^2(F_{\rm av})
=
\begin{pmatrix}
0&0&1\\
0&1&2\\
1&1&1
\end{pmatrix},
\qquad
\operatorname{Spec}\bigl(\operatorname{Sym}^2(F_{\rm av})\bigr)
=\{\varphi^2,-1,\varphi^{-2}\}.
\]
The eigenvalue \(\varphi^{-2}\) is the positive stable mode. On closed
returns, the scale cocycle of \cref{sec:cociclo-parry-escala} gives
\[
M_n=\varphi^{-2n}={\bigl(\chi_P(\gamma)\bigr)}^{-2}.
\]

\begin{proposition}[Functional content of \(\varphi\)]
\label{prop:autoescala-phi-rev6}
The function of \(\varphi\) in HMT is the self-scaling of the quotient of modes:
growth by the Perron root \(\varphi\), cylindrical conservation by
\(\chi_P\) and stable second-order return by \(\varphi^{-2}\).
These three expressions belong to one and the same operator induced by the
TPK calendar.
\end{proposition}

\begin{proof}
The incidence \(F_{\rm av}\) descends from the block
\((\Sigma,\Pi,\Pi)\) by
\cref{thm:descenso-necesario-fav}. Its characteristic polynomial produces the
Perron root, and the multiplicativity of \(\chi_P\) is obtained by
telescopy. The quadratic action separates the three indicated eigenvalues; the
contracted positive mode is \(\varphi^{-2}\).
\end{proof}

\section{Oriented correspondence between closure and self-scale}

The matrix \(F_{\rm av}\) generates \(\varphi^{-2}\), but cannot by itself
generate \(\pi\): every rational function of its spectrum belongs to
\(\mathbb Q(\sqrt5)\). Transcendental closure comes from the coinductive regional
reader. This separation of origins is preserved by the diagonal operator
\[
 D_\pi=\pi_{\rm HMT}P_{\rm ar}+P_{\rm vol},
\qquad
 P_{\rm ar}=
 \begin{pmatrix}1&0\\0&0\end{pmatrix},
\quad
 P_{\rm vol}=
 \begin{pmatrix}0&0\\0&1\end{pmatrix}.
\]
The marked return of length \(n\) is
\[
\Theta_n
=\pi_{\rm HMT}\frac{F_{n-1}}{F_{n+1}},
\qquad
\lim_{n\to\infty}\Theta_n
=\frac{\pi_{\rm HMT}}{\varphi^2}.
\]
The dynamical proof, the exact witness of length \(108\), and the reader
of nested cylinders reside in
\cref{sec:retorno-marcado-pi-phi2}. The mark of \(\pi\) is applied only once
to the return: it is not inserted into every transition of \(F_{\rm av}\).

The same reason has two orientations. Let
\[
\mathscr R(x,y)=\frac{x}{y^2},
\qquad
\mathsf J(x,y)=(y,x).
\]
Then \(\mathsf J^2=I\) and
\[
\boxed{
\mathscr R(\pi,\varphi)=\frac{\pi}{\varphi^2},
\qquad
(\mathscr R\circ\mathsf J)(\pi,\varphi)
=\frac{\varphi}{\pi^2}.}
\]
Moreover,
\[
\frac{\pi/\varphi^2}{\varphi/\pi^2}
={\left(\frac{\pi}{\varphi}\right)}^3,
\qquad
\frac{\varphi}{\pi^2}
=\frac{1}{\varphi^3{\bigl(\pi/\varphi^2\bigr)}^2}.
\]
The first orientation reads areal closure over volumetric memory; the second, interior self-scaling over quadratic closure. Section~\ref{sec:espejo-pi-phi-electron-hmt} proves this involution and its subsequent appearance in the electron exponent. This correspondence does not identify \(\pi\) with \(\varphi\): it preserves their order and makes explicit what changes when that order is reversed.

\section{Arbitrary precision of monodromic emission}
\label{sec:precision-arbitraria-monodromia-rev6}

\begin{corollary}[Emission of any finite prefix]
\label{cor:precision-arbitraria-monodromia-rev6}
Under the hypotheses of
\cref{thm:generacion-supervivencia-pi-e-phi-rev6}, for every \(M\geq1\) the
EF--\(G_9\) iteration determines the first
\(M\) decimal digits of each of the three channels. In particular, \(M=100\)
and \(M=1000\) are finite windows of the same monodromic prolongation and do not
require the introduction of an external table of digits.
\end{corollary}

\begin{proof}
Let \(K(n)\) be the number of triads published after \(n\) ternary
blocks. The section supplied by the theorem has levels of every
depth and \(K(n)\to\infty\). Given \(M\),
choose \(n\) with \(3K(n)\geq M\). The
\cref{thm:generacion-supervivencia-pi-e-phi-rev6} uniquely determines
the truncation of depth \(n\), and its decimal cylinder fixes the
first \(3K(n)\) digits. Taking the first \(M\) concludes the argument.
\end{proof}

The difference of functions is fixed within a single construction: \(\pi\) closes distinct regional provenances in a common evaluation; \(e\) propagates a visible output while renewing the memory that transports it; \(\varphi\) governs the self-scaling of the quotient of modes. Monodromy explains how those functions return to the same phase and pass through the mirror \(\pi/e\). Under the coinductive hypotheses of the prolongation theorem, they publish new blocks indefinitely without resetting their history.

% El resultado de realizabilidad de una cinta prescrita se conserva como
% material auxiliar, pero no forma parte del argumento generativo rector.
\ifdefined\HMTIncludeTargetFedRay
\section{Inductive prolongation by rays: existence, compatibility, and domain}
\label{sec:rayo-coinductivo-rev6}

The passage from every finite prefix to an infinite route uses a theorem of
finitely branched compactness. Its object is not the selector of the catalogue,
but the realization of a channel whose expansion has already been fixed.

Let
\[
\mathcal S
=\mathbb Z_9^2
\times\{+,\times\}
\times\{N,E,S,O\},
\qquad
|\mathcal S|=648,
\]
the finite state space of the route model. The control alphabet is
\(\mathcal A=\{N,E,S,O\}\), and the cost of a route counts turns:
\[
c(a_{t-1},a_t)=\mathbf1_{\{a_t\neq a_{t-1}\}}.
\]
Fix a channel \(x\in\{\pi,e,\varphi\}\) and its ternary expansion, and let
\(W_m(x)\) be a canonical horizon-\(27m\) route that minimizes this cost
among the routes realizing the first \(9m\) trits of \(x\).

\begin{theorem}[Limit Ray of a Fixed Expansion]
\label{thm:rayo-coinductivo-condicionado-rev6}
For each fixed expansion \(x\in\{\pi,e,\varphi\}\), there exists a route
\[
\mathcal D_\infty(x)\in\mathcal A^{\mathbb N}
\]
such that:
\begin{enumerate}
\item for every \(L\geq1\), the length-\(L\) prefix of
\(\mathcal D_\infty(x)\) coincides with the prefix of \(W_m(x)\) for
infinitely many values of \(m\);
\item \(\mathcal D_\infty(x)\) realizes exactly the infinite ternary
expansion of \(x\).
\end{enumerate}
Among the branches with the first property, the lexicographically minimal one is taken.
\end{theorem}

\begin{proof}
For each \(L\), consider the set \(T_L(x)\) of prefixes of length
\(L\) that occur at the beginning of \(W_m(x)\) for infinitely many values of
\(m\). It is nonempty: there are only finitely many prefixes of length \(L\) and
infinitely many horizons. The sets \(T_L(x)\), ordered by extension,
form a finitely branching tree. König's lemma yields an infinite
branch; the least in lexicographic order defines
\(\mathcal D_\infty(x)\).

Once \(n\) has been fixed, take \(L\) sufficient to determine the first \(n\)
filtering trits. That prefix appears in infinitely many \(W_m(x)\), and each of
them realizes by definition the first \(n\) trits of \(x\). The same
prefix of \(\mathcal D_\infty(x)\) realizes them. Since \(n\) is arbitrary,
the route realizes the complete expansion.
\end{proof}

Now let
\[
\mathcal R(x)=\{729x\},
\]
where \(\{\cdot\}\) denotes fractional part. Multiplication by
\(729=3^6\) shifts six trits. In the model's calendar, each filtration trit
occupies three elementary iterations, so that six trits
are equivalent to eighteen elementary iterations.

\begin{corollary}[Covariance by the self-map \(729\)]
\label{cor:covarianza-rayo-729-rev6}
Let \(\sigma\) be the left shift of a route. For every
\(r\geq0\),
\[
\sigma^{18r}\mathcal D_\infty(x)
\]
performs exactly the channel \(\mathcal R^r(x)\).
\end{corollary}

\begin{proof}
Finite covariance identifies the suffix that begins after \(18r\)
elementary iterations
with the displacement of \(6r\) trits. The
\cref{thm:rayo-coinductivo-condicionado-rev6} transports that identity to
all prefixes of the limit ray.
\end{proof}

The corollary expresses an iterated holography: one and the same finite cell,
displaced by \(\mathcal R\), realizes successive strata of the expansion.
It does not turn \(\mathcal R\) into a selector of \(\pi,e,\varphi\), because
\(x\) already intervenes when defining the sets of admissible routes
\(W_m(x)\). Its exact function is subsequent: once the channel has been distinguished,
it constructs a compatible ray and makes its covariance across scales visible.

\section{Orbital selector composition, the numeric cylinder, and Archimedean evaluation}

The three final operations form a causal diagram, not three versions
of the same theorem:
\[
\text{internal state}
\xrightarrow{\,\mathfrak S_{\rm cat}\,}
w_x
\xrightarrow{\ \text{enriched reader}\ }
\text{expansion }x
\xrightarrow{\,\mathcal D_\infty(x)\,}
\text{ray that realizes it}.
\]
The orbital selector \(\mathfrak S_{\rm cat}\) distinguishes the three words
by catalogue invariants and an explicit orienting gauge. The enriched reader
prolongs word, cylinder, and memory. The theorem of rays
then constructs an optimal recurrent infinite realization of the expansion
already fixed.

\begin{proposition}[Separation of force among the three layers]
\label{prop:tipado-biografias-rayos-selector-rev6}
The following assertions hold simultaneously.
\begin{enumerate}
\item The five regions and the minimal closures of \(\pi\), the local square
and the memory advance without state periodicity of \(e\), and the self-scale
\(F_{\rm av}\)--\(\varphi\)--\(\varphi^{-2}\) are exact results in
their declared domains.
\item The ray \(\mathcal D_\infty(x)\) is an exact coinductive realization
of a fixed expansion.
\item The rayo not is a selector autonomo, because the expansion \(x\) appears
in the definition of \(W_m(x)\).
\end{enumerate}
The third assertion limits only the selector use of the theorem of
rays; it does not modify the first two.
\end{proposition}

\begin{proof}
The results of the first point are proved, respectively, in
\cref{chap:regiones,chap:cinco-ciclos-pi,thm:cuadrado-local-e,%
thm:descenso-necesario-fav,cor:parry-envolvente-fav}. The second point is
the \cref{thm:rayo-coinductivo-condicionado-rev6}. For the third, it suffices to
observe that the minimization set defining \(W_m(x)\) requires
realizing the first \(9m\) trits of \(x\); therefore, that same route cannot be used
to deduce which \(x\) should have been selected.
\end{proof}

The difference of functions is thereby definitively fixed. \(\pi\) closes
distinct regional histories into a single evaluation; \(e\) propagates a
visible output while renewing the memory transporting it; \(\varphi\)
governs the scale change of the mode quotient. The correspondence
\(\pi/\varphi^2\leftrightarrow\varphi/\pi^2\) coordinates closure and
self-growth without conflating them. The coinductive ray explains how an
already distinguished expansion is prolonged indefinitely. This is the biography:
not only what number appears, but what structure selects it, what memory it
preserves, and what function it fulfills within the system.
\fi
